\title{X-LoRA: Mixture of Low-Rank Adapter Experts, a Flexible Framework for Large Language Models with Applications in Protein Mechanics and Molecular Design}

\begin{abstract}
We present a mixture of expert approach to develop fine-tuned large language models employing a deep, layer-wise token-level technique based on low-rank adaptation (LoRA). Beginning with a collection of pre-trained LoRA adapters, our gating mechanism leverages hidden states to dynamically blend adapted layers, enabling the resulting X-LoRA model to harness diverse capabilities and generate novel deep layer-wise combinations for task resolution. This framework draws inspiration from biological concepts of universality and diversity, where neural network components are reused across various hierarchical forms. Consequently, the X-LoRA model can be readily applied to any existing large language model (LLM) without requiring modifications to the underlying architecture. We design a specialized X-LoRA model offering scientific functionalities including forward and inverse analysis tasks alongside enhanced reasoning abilities, particularly focused on biomaterial analysis, protein mechanics, and design. The significance of this work lies in providing accessible, expandable, and adaptable models with rich domain expertise and the ability to integrate knowledge from multiple fields. Incorporating experts in biology, mathematics, reasoning, bio-inspired materials, mechanics and materials science, chemistry, protein biophysics, mechanics, and quantum mechanics-based molecular properties, we perform a series of physics-centered case studies. These include assessments of knowledge recall, protein mechanics forward and inverse problems, protein design, adversarial agentic modeling with ontological knowledge graph construction, and molecular design. The model not only generates quantitative predictions of nanomechanical properties of proteins and quantum mechanical molecular characteristics but also performs reasoning on these outcomes, accurately identifying plausible mechanisms explaining distinct molecular phenomena.
\end{abstract}

\section{Introduction}
Large language models (LLMs)~\cite{Vaswani2017AttentionNeedc,Touvron2023LlamaModels,OpenAI2023GPT-4Report,Chowdhery2022PaLM:Pathways,Jiang2023Mistral7B,Gunasekar2023TextbooksNeed} have become increasingly popular, including the creation of specialized models that excel in specific task types, reasoning processes, or scientific fields~\cite{Bubeck2023SparksGPT-4,Buehler2023MechGPTModalities,Nejjar2023LLMsAnalysis,Buehler2023GenerativeDesign,Luu2023BioinspiredLLM:Materials,Luu2023GenerativeSolvents,Buehler2023MeLMProblemsb,Ge2023OpenAGI:Experts}. However, training these models can be expensive, particularly when a wide range of capabilities is required. Approaches like low-rank adapters (LoRA)~\cite{hu2021lora} have been introduced as a more resource-efficient alternative, though these adaptations generally target more specialized domains. The key idea behind LoRA is to use low-rank matrices added to the original full-scale matrix and treat these low-rank matrices as the sole trainable parts of the model. Because only the adapter layers are updated, such models tend to retain the knowledge gained during the base model’s pre-training while becoming better suited for particular tasks, all while reducing computational costs. Given the efficiency of LoRA adapter training, it is feasible to build a broad array of specialized models through this technique.

Although LoRA adapters offer an effective means to develop specialized model capabilities, the challenge remains on how to combine multiple such models into enhanced models that provide an integrated and potentially superior range of abilities. Indeed, the idea of combining LLMs has attracted growing research interest~\cite{Kim2023SOLARUp-Scaling}. Experimental methods, such as using spherically interpolated model parameters (SLERP) and other techniques~\cite{Arcee-ai/mergekit:Models.}, have demonstrated promising outcomes and indicate potential for combining models to unlock new functionalities. Typically, these mixed models are generated {\it a priori} following a defined algorithm, and although they have shown encouraging results on certain benchmarks, further investigation is required to systematically understand their construction and the emergence of particular features.

We propose that dynamically mixing parameters could serve as a more universal method, enabling the combined model to discover novel parameter combinations during inference.

Similar methodologies have employed mixture of experts (MoE) frameworks~\cite{Jacobs1991AdaptiveExperts,Hampshire1992TheRecognition,Jordan1994HierarchicalAlgorithm}, including the recently introduced Mixtral LLM~\cite{Jiang2024MixtralExperts}. Nonetheless, these approaches can be computationally intensive, requiring substantial resources even at inference time. In contrast, we present a computationally efficient mixture-of-experts technique that utilizes a collection of distinct LoRA adapters, each of which can be independently trained with ease. Our method embodies a high-resolution and adaptable mechanism, drawing inspiration from a biological model that reutilizes neural network components to build a hierarchy of interacting models, spanning from foundational models to agentic systems (see Fig.~\ref{fig:hierarchical_concept}). To quantitatively demonstrate its effectiveness, our study targets the application of this approach to science-oriented LLMs, specifically addressing protein mechanics challenges within the broader context of bio-inspired materials analysis and design.

We introduce an algorithm that enables dynamic mixing of adapters at the token-level state, allowing for fine-grained combination of adapters down to the level of individual layers. Termed X-LoRA, this method grants access to a rich variety of functionalities essential for scientific applications involving multimodal language modeling~\cite{Hu2023DeepScience,Buehler2023AMetamaterials,Buehler2022ProteinNetworks}.

\subsection{Fundamental concepts of X-LoRA}

We begin by briefly revisiting the principles underpinning LoRA adapters. The core idea behind low-rank adaptation (LoRA)~\cite{hu2021lora} posits that weight updates lie in a low “intrinsic dimension.” Exploiting this, the approach fixes the original weights $W_0 \in \mathbb{R}^{d \times k}$ and restricts the updates to a low-rank decomposition

\begin{equation}
W_0 + \Delta W_0 = W_0 + BA 
\end{equation}

where $B \in \mathbb{R}^{d \times r}$ and $A \in \mathbb{R}^{r \times d}$, with the rank $r \ll {\rm min}(d,k)$. When training with LoRA~\cite{hu2021lora}, the pretrained weights $W_0$ remain fixed, and only the matrices $A$ and $B$ contain parameters that are updated. The forward pass for a LoRA layer can be written as:

\begin{equation}
    h = W_0x + \Delta W_0x = W_0x + BAx
\end{equation}

In practical usage, LoRA adapters include a fixed scalar weight $\alpha$, which is applied to the decomposition matrices, resulting in:

\begin{equation}
    h = W_0x + BAx \cdot \alpha
\end{equation}

Building on this concept, X-LoRA introduces a set of LoRA adapters, each endowed with distinct capabilities, as demonstrated in previous scientific LLM studies (e.g., \cite{Buehler2023GenerativeDesign}). Instead of statically mixing or integrating models, we propose a dynamic gating mechanism that adjusts each LoRA adapter's scaling at the level of individual tokens and layers, allowing fine-grained combination deep within the model. Specifically, for a LoRA adapter $i$, the original scaling factor $\alpha_i$ is modified by a factor $\lambda_i$ to produce the updated scaling $\alpha^*_i$:

\begin{equation}
    \alpha^*_i = \alpha_i \cdot \lambda_i
\end{equation}

The scaling factor $\lambda_i$ is generated by an X-LoRA scaling head, which leverages the model’s hidden states and constitutes the trainable element in the X-LoRA framework. Since the scaling head can process complex representations, it can be efficiently implemented as a simple feed-forward neural network with one or a few layers (users can customize the architecture of this network in the X-LoRA implementation).

This scaling computation is performed separately for each token in the input sequence, thus providing fine granularity by scaling each LoRA adapter operating at specific layers within the LLM.

From here on, we denote this method of gating the various adapter components via scaling as simply “scaling.” Additional details about this method are presented in the Materials and Methods section.

\subsection{Paper outline}

The structure of this paper is organized as follows. Initially, we describe the methodology and training approach, including the pertinent datasets, which lead to the creation of an X-LoRA model possessing specialized capabilities in the physical sciences, with a particular emphasis on biomaterials such as proteins. Subsequently, we showcase a series of experiments where the X-LoRA model is applied to diverse tasks, including question answering, conversational and agent-based modeling, protein design and analysis, among others. We perform an in-depth examination of the scaling behaviors and validate the methodology by comparing it with molecular modeling and other physical data and techniques.

\section{Results and discussion}

The development of the X-LoRA model involves the following stages:

\begin{enumerate}
    \item Training the foundational base LLM (completed in previous work) \item Training individual adapters to build expertise across several domains (these represent distinct or overlapping skill and knowledge areas) \item Training the integrated X-LoRA model
\end{enumerate}

Our experimental work begins with the training of multiple LoRA adapters. We create a set of nine adapters, each fine-tuned with specialized expertise, based on the Zephyr-7B-$\beta$ model~\cite{HuggingFaceH4/zephyr-7b-betaFace}, which itself was developed upon the Mistral-7B model\cite{Jiang2023Mistral7B}:

\begin{enumerate}
    \item Bioinspired materials \item Chain-of-thought (CoT) and reasoning \item Chemistry \item Mathematics \item Physics \item Biology \item Mechanics and materials \item Platypus: Logic and reasoning \item Protein mechanics tasks (generative sequence-to-property and inverse capabilities)
\end{enumerate}

These chosen areas of focus and expertise were selected because our ultimate goal is to develop a STEM-oriented LLM with an emphasis on materials science. Nonetheless, we hypothesize that our method can be broadly applied across various fields, particularly to establish modalities that facilitate connections between different domains of knowledge~\cite{Buehler2023MechGPTModalities,Buehler2023MeLMProblemsb}.

Once the LoRA adapters are generated, we proceed to train the X-LoRA model. For this training, we utilize a subset of the combined training samples originally used to train the individual LoRA adapters (further details are provided in Materials and Methods, including comprehensive descriptions of the dataset, implementation, and related aspects).

The X-LoRA system serves as an experimental framework to evaluate its effectiveness under various scenarios. We perform a series of tests, which include tackling complex question answering, employing the generative model for protein design, conducting analysis and reasoning tasks, as well as developing more advanced interaction strategies via adversarial agentic modeling. Throughout our evaluation, we compare the performance of the base model with that of the X-LoRA model, and monitor how the scaling head engages different layers from the various experts during the generation or formulation of answers.

\subsection{Question answering and observed X-LoRA layer-wise scaling weights}

Initially, we analyze the distinctions between the X-LoRA model and the foundational base model. We present both models with the following question:

\begin{quote}
\small\texttt{What material behavior (stiffening vs softening) causes supersonic crack speeds in mode II cracks? Think step by step.}
\end{quote}

The base model replies:

\begin{quote}
\small\texttt{Supersonic crack speeds, also known as Rayleigh waves, occur in mode II cracks due to the stiffening material behavior of the material near the crack tip. In mode II fracture, the loading is shear-dominated and the crack propagates perpendicular to the direction of applied stress. Here's a step-by-step explanation:

As the crack tip nears the material, the stresses around it escalate swiftly. This generates an area of pronounced stress concentration known as the stress singularity.

Within this stress singularity, the material undergoes elastic deformation that results in increased stiffness, or stiffening behavior. This stiffening arises because the material adjacent to the crack tip becomes more restricted as the crack advances [...]}
\end{quote}

Drawing from our own expertise~\cite{Buehler2008AtomisticFailure}, this explanation is not only inaccurate but also unnecessarily verbose. By contrast, the X-LoRA model provides the following response:

\begin{quote}
\small\texttt{The material behavior responsible for supersonic crack propagation in mode II cracks is stiffening.}
\end{quote}

This response is accurate, concise, and clear~\cite{Buehler2008AtomisticFailure}.

The contrast between the base model and the X-LoRA model becomes even more pronounced in specialized protein analysis and design tasks. For example, the task presented is:

\begin{quote}
\small\texttt{CalculateForceHistory< M S K S N T Y R M L V L L E D D T K I N K E D E K F L K G K P G K M H E F V D E L I L P F N V D E L D E L N T W F D K F D A E I C I P Y V E I L K E S G M K > }
\end{quote}

The response from the base foundation model is:

\begin{quote}
\small\texttt{I'm sorry but the input provided is not a valid mathematical expression. It appears to be a sequence of characters that do not form a calculable function. Please provide a valid mathematical expression for me to calculate.

However, I can tell you that the sequence of characters you provided is actually a series of function and variable names that may be part of a larger program or script in a specific programming language. Without context, it's impossible to determine the intended calculation.

If you provide more information, such as the programming language and the context in which this expression appears, I may be able to provide more insight into its meaning and possible calculations. }
\end{quote}

The X-LoRA model replies as follows, successfully accomplishing the design objective of producing the force-deformation curve related to this protein:

\begin{quote}
\small\texttt{[0.013,0.066,0.078,0.045,0.072,0.126,0.127,0.107,0.119, 
0.074,0.106,0.098,0.103,0.128,
0.077, 0.069,0.131,0.108,0.144,0.174, 0.166,0.183,0.185,0.217,0.175,0.230,0.187,0.224, 0.267,0.301,0.315]}
\end{quote}

Figure~\ref{fig:QA_weights_sample_0} presents outcomes from this question answering along with the observed X-LoRA scaling weights. The two queries—one centered on mechanics and materials (Fig.~\ref{fig:QA_weights_sample_0}a) and the other on protein mechanics (Fig.~\ref{fig:QA_weights_sample_0}b)—exhibit distinct patterns. Notably, we observe that in both cases, a complex arrangement of scaling values is employed, indicating that the X-LoRA model leverages heterogeneous mixing of different adapters across layers. Furthermore, the heatmap in Fig.~\ref{fig:QA_weights_sample_0}a reveals a wider spread of activation over many experts and layers, featuring some regions with strong activation. Conversely, the heatmap in Fig.~\ref{fig:QA_weights_sample_0}b displays a more focused activation pattern, with fewer areas of high intensity. We interpret this as possibly reflecting that the gating function on the right is more selective in which experts it activates.

In Fig.~\ref{fig:QA_weights_sample_0}, the areas resembling bright lines indicate pathways of high activation, signifying that the scaling function preferentially activates certain experts at specific layers. The observed patterns imply that the importance is not uniformly distributed among all experts and layers; rather, it tends to be sparse and selective. Such sparsity is a common feature in other mixtures-of-experts models, where only a subset of experts is selected for each input to focus on distinct parts of the problem space.

By comparing these maps across different tasks, both the sparsity and the specific patterns could be crucial for gaining deeper insight into which experts the model depends on for various input types or at different stages of processing within the neural network (we will later discuss how these maps evolve dynamically during the resolution of complex tasks).

To deepen the analysis, we now consider a wider range of questions to investigate how the X-LoRA model engages different experts in response to a variety of prompts. Our focus is particularly on whether, and if so how, the X-LoRA model blends different adapters when dealing with questions that span overlapping skills or domains. Fig.~\ref{fig:QA_weights} displays results from question answering along with the observed X-LoRA scaling weights, corresponding to the queries outlined in Table~\ref{tab:table_qa}. In Table~\ref{tab:table_qa}, we present both the responses from X-LoRA and the baseline model, Zephyr-7B-$\beta$. The differences between these two models are quite pronounced.

To illustrate this clearly, we highlight incorrect or questionable responses with a \redhighlight{red highlight}. This comparison demonstrates that X-LoRA delivers significantly higher accuracy and more succinct answers.

As anticipated, we notice intricate mixing of adapters and frequent activation of multiple dominant LoRA experts. For example, Fig.~\ref{fig:QA_weights}d shows that a question regarding the mechanics of pine cones triggers strong activation of the mechanics/materials adapter together with the CoT/reasoning and bioinspired adapters. Although X-LoRA selects the correct letter, the answer is not fully accurate since pine cones open under dry, not wet, conditions. However, a follow-up question probing whether release occurs in dry or wet conditions produces the correct answer.

For a protein-related query, Fig.~\ref{fig:QA_weights}f shows adapter mixing in response to a question about the mechanics of Bombyx mori silk, activating a combination of the bioinspired, mechanics/materials, and biology adapters. In cases involving clearly specialized tasks such as protein design (Fig.~\ref{fig:QA_weights}a), we observe predominantly singular activation of a single expert.

A more detailed examination of the heatmap of scalings shown in Fig.~\ref{fig:QA_weights}a reveals that beyond the clearly significant importance of the protein mechanics expert, there exist several bands where specific experts receive substantial weighting across multiple layers. This phenomenon is especially prominent for the mechanics/materials and reasoning experts. These bands are discontinuous and exhibit a fluctuating pattern, which may indicate particular conditions or features that these experts manage effectively. Notably, the reasoning expert (positioned immediately to the left of the protein mechanics expert) is distinguished by multiple layers where it attains the highest activation, particularly in the higher layers (approximately layers 100 to 140). We hypothesize that this could suggest a crucial role in the model’s final decision-making or output, though further investigation is required to confirm this. Additionally, there is a more localized yet strong activation for the physics and biology experts at specific layers, which might reflect a specialization of these experts for certain features or representations within those layers.

Moreover, Fig.~\ref{fig:QA_token_history} illustrates a token-level history based on the aggregate scalings summed over all layers. This analysis presents token histories for two cases: in Fig.~\ref{fig:QA_token_history}a, a task associated with pine cone seed release, and in Fig.~\ref{fig:QA_token_history}b, the outcome of a sequence of protein mechanics tasks. The protein mechanics task involved a series of prompt-answer exchanges, beginning with two property calculations followed by a generative task aimed at designing a novel protein towards the end. Although expert usage remains generally stable, there are noticeable shifts in the involvement of the LoRA experts throughout the process. In particular, the protein mechanics expert is prominently engaged during the final stages of generation when a new protein design is being created.

Table~\ref{tab:table_qa} presents a variety of use cases across multiple domains along with a performance comparison against the base model alone. We explore several of these cases in greater detail.

For example, we present a question that combines elements of protein science and biology with logical reasoning, focusing on protein expression patterns. X-LoRA provides the correct answer. The base model’s response, however, contains some errors and logical flaws when evaluating the expression of marker proteins in each cell type. The base model accurately notes that fibroblasts express Protein B and do not express Protein A due to the mutual exclusivity between Protein A and Protein B. Nonetheless, it fails to explicitly state that fibroblasts also express Protein C. Since there is no rule preventing Protein C expression alongside Protein B, and fibroblasts have no restriction regarding Protein C, it follows logically that they express Protein C as well. Regarding neurons, the base model correctly begins by asserting that neurons express Protein A and not Protein B because of the mutual exclusivity rule. However, its discussion about neurons potentially expressing Protein C or both Protein C and Protein B is unclear and contains an error. Given that neurons express Protein A and not Protein B, the only remaining possibility, according to the rules, is that they express Protein C. There is no condition under which neurons could express Protein B, contrary to the initial analysis. Thus, neurons express Protein A and Protein C, but not Protein B. For epithelial cells, the base model rightly states that these cells do not express Protein A but do express both Protein B and Protein C. However, the explanation is overly complicated and includes incorrect claims regarding the mutual exclusivity of Protein A and Protein C, which was not part of the original information. The straightforward reasoning is that epithelial cells do not express Protein A but must express the other two proteins, Protein B and Protein C, since the initial instructions only prohibit co-expression of Protein A and Protein B.

X-LoRA, by contrast, accurately determines the marker proteins expressed by each cell type based on the given observations and logical reasoning. Regarding fibroblasts, X-LoRA correctly infers that fibroblasts cannot express Protein A because the presence of Protein A rules out the expression of Protein B. Since fibroblasts do express Protein B, and considering that each cell type expresses a unique combination of marker proteins (the assumption that this combination includes exactly three proteins was incorrect and not stated in the original problem, which only indicated uniqueness), the conclusion that fibroblasts also express Protein C is valid based on the information provided (except for the misunderstanding about the number of proteins). For neurons, the model rightly identifies that they express Protein A and Protein C. This reasoning is logical, given that Protein B cannot coexist with Protein A, and neurons express Protein A plus another marker protein that is not Protein B, thus it must be Protein C. Lastly, for epithelial cells, X-LoRA correctly concludes that they express Protein B and Protein C since epithelial cells do not express Protein A and, by elimination, express two marker proteins which must be Protein B and Protein C.

In the question regarding pine cone release under humid versus dry conditions, the base model offers a confusing and factually incorrect response, whereas X-LoRA provides a clear and precise explanation. Other related cases include the question about the exceptional mechanical properties of Bombyx mori silk, where the base model mistakenly claims that silk fibers conduct electricity. Additionally, the base model fails at several domain-specific tasks, such as the question concerning the role of hyperelasticity in maximum fracture speeds, where X-LoRA delivers a succinct and accurate answer, but the base model does not respond appropriately. In the question about the strong adhesion of gecko feet, X-LoRA correctly highlights the importance of setae and nanoscale effects, while the base model incorrectly attributes the adhesion to special-purpose glue proteins. These examples, along with others detailed in the table and throughout the paper, provide compelling evidence for the superior performance of X-LoRA compared to the base model.

Figure~\ref{fig:perf_comp_bioinspiredexam} presents a comparison of knowledge recall assessment experiments, utilizing the bio-inspired knowledge test set introduced in \cite{Luu2023BioinspiredLLM:Materials}. It is evident that X-LoRA demonstrates the highest performance, despite being a substantially smaller model than either the BioinspiredLLM or the Llama-BioLLM models (7B parameters for X-LoRA compared to 13B parameters for the others). The figure also includes comparisons with several other recently released LLMs, specifically Orca-13B and Llama-13b-chat.

Additionally, the plot compares X-LoRA to its base model, Zephyr-7B-$\beta$, highlighting that the X-LoRA model achieves significantly enhanced performance. These results suggest that strong outcomes can be attained with X-LoRA even when smaller base models are employed.

Figure~\ref{fig:image_synth} illustrates an application in image synthesis, where X-LoRA is utilized to generate a prompt that can be fed into other models, such as DALL-E 3, as demonstrated here.

\begin{quote}
\small\texttt{Create a prompt that I can use for image generation that accurately describes the microstructure of a hierarchical composite.

Explicitly describe the geometry.}
\end{quote}

The comparison shown in Figure~\ref{fig:image_synth}(b)-(c) clearly indicates that the prompt produced by X-LoRA is more concise, which leads to a significantly more accurate image generated by DALL-E 3~\cite{DALLECard}. Notably, the prompt generated by Zephyr-7B-$\beta$ fails to precisely follow the given instruction and introduces various additional design elements such as fibers and nanoparticles. These elements were not mentioned in the original prompt, which specifically focused on the microstructure of a hierarchical composite. The X-LoRA prompt is both shorter and more targeted, resulting in a markedly improved outcome.

\subsection{Generative solutions to protein analysis}

We now turn our attention to particular protein generative and analytical tasks. These functionalities emerge in the X-LoRA model through the protein mechanics adapter and are further enhanced by integrating the model's knowledge of biology, bio-inspired materials, reasoning, and logical deduction. This section focuses exclusively on analyzing protein tasks, evaluating the quantitative performance of this capability.

Fig.~\ref{fig:protein_force_history_prediction} presents outcomes derived from the protein task, which involves predicting force-deformation behaviors based on sequences. Fig.~\ref{fig:protein_force_energy_prediction} illustrates the overall accuracy in forecasting unfolding force and unfolding energy from the sequence, comparing actual values with predictions (yielding an $R_2=0.85$). In general, the model exhibits excellent performance and, as will be examined in the following section, also performs very effectively for the inverse task of design and cycle consistent evaluation\cite{Lu2023GenerativePropertiesb}. Nevertheless, unlike previous studies that employed highly specialized GPT-type models to address these tasks, our X-LoRA model integrates these specialized capabilities alongside a range of other deep specializations introduced through various fine-tuned adaptations.  

\subsection{Protein design and analysis}

We now broaden the scope of test cases to encompass more integrative applications involving multiple capability areas. Specifically, the experiments described in this section demonstrate how the X-LoRA model leverages a unique combination of skills and how agentic modeling can yield even more complex and comprehensive responses. The discussion begins with a protein design task, in which we request the model to design a novel protein that achieves a specific target force-deformation response. These force-deformation responses have been previously examined~\cite{Ni2023ForceGen:Model} and correspond to measurements of force versus deformation as the protein is stretched at both ends, reflecting a standard protein unfolding experiment~\cite{Lu1998UnfoldingSimulation.}.

Fig.~\ref{fig:design_protein} illustrates the application of the generative protein task. We create a protein exhibiting a targeted force-deformation response (using training data and boundary conditions from the original molecular dynamics (MD) simulations of protein pulling, see~\cite{Ni2023ForceGen:Model}), then evaluate the predicted sequence to verify if it achieves the intended target characteristics (Fig.~\ref{fig:design_protein}a). This cycle-consistent evaluation is crucial for determining whether the inverse design aligns with the forward prediction.

Fig.~\ref{fig:design_protein}b displays the sequence predicted: \texttt{MSKSNTYRMLVLLEDDTKINKEDEKFLKGKPGKMHEFVDELILPFNVDELDELNTWFDKFDAEICIPYVEILKESGMK}. Additionally, Fig.~\ref{fig:design_protein}c presents a folded structural model of the protein, generated using AlphaFold 2~\cite{Jumper2021HighlyAlphaFoldb}. To explore the relationship of the designed protein to known sequences, Fig.~\ref{fig:design_protein}d shows a Basic Local Alignment Search Tool (BLAST) Tree~\cite{Altschul1990BasicTool} analysis. Notably, the designed protein demonstrates certain connections to hypothetical proteins examined in previous studies but which have neither been identified nor structurally resolved.

By employing agentic modeling and integrating multiple X-LoRA models as components within a multi-agent framework, we conduct further analysis of the predicted sequence using our X-LoRA model. This approach leverages our model’s enhanced capacity not only to comprehend protein mechanics but also to combine this understanding with knowledge from other domains. The prediction obtained above serves as input, and we request a more detailed analysis along with design recommendations for sequence variants. We begin with the following prompt:

\begin{quote}
\small\texttt{I have identified this amino acid sequence: MSKSNTYRMLVLLEDDTKINKEDEKFLKGKPGKMHEFVDEL ILPFNVDELDELNTWFDKFDAEICIPYVEILKESGMK.

How can I use it?

Provide logical steps, and discuss the engineering aspects of the problem specifically.}
\end{quote}

Multiple agentic interactions are utilized to formulate the solution, employing adversarial strategies (see Materials and Methods for details) to encourage the model to delve into deeper and more complex aspects of the discussed concepts. A summary of this dialogue, as identified by the X-LoRA model, is provided below (for the complete conversation, see Supplementary Material, SI Conv Example 1):

\begin{quote}
\small\texttt{The dialogue explores the possible uses of a particular amino acid sequence in the field of protein engineering. This sequence features a combination of hydrophobic and hydrophilic amino acids, indicating its potential to form a stable conformation and interact with other molecules. To design a protein containing this sequence, computational approaches can be employed to predict its secondary and tertiary structures; however, challenges arise due to the complexity of protein folding and the high computational demand. Enhancements can be achieved by creating more efficient algorithms and integrating novel experimental data. Ensuring the accuracy and practical relevance of computational models requires thorough validation. Methods such as cross-validation, regularization, early stopping, ensemble techniques, and transfer learning are useful to avoid overfitting and underfitting. Additionally, model interpretation and explanation can provide valuable insights for real-world applications. }
\end{quote}

We also ask the model to develop a list of key points, resulting in:

\begin{quote}
\small\texttt{\begin{itemize}
    \item The amino acid sequence MSKSNTYRMLVLLEDDTKINKEDEKFLKGKPGKMHEFVDELI LPFNVDELDELNTWFDKFDAEICIPYVEILKESGMK holds promise for protein engineering owing to its blend of hydrophobic and hydrophilic residues. Computational techniques can assist in predicting the protein’s secondary and tertiary structures, though limitations stem from the intricate nature of protein folding and computational expense.
    \item Validating computational models is essential to guarantee their precision and applicability in practical scenarios. Strategies like cross-validation, regularization, early stopping, ensemble methods, and transfer learning can help mitigate overfitting and underfitting.
    \item Understanding and clarifying models can further aid in guiding practical implementations.
\end{itemize}
}
\end{quote}

Finally, the key takeaway is:

\begin{quote}
\small\texttt{The key insight is that the amino acid sequence MSKSNTYRMLVLLEDDTKINKEDEKFLKGKPGKMHEFVDELILPFNVDELDELNTWFDKFDAEICIPYVEILKESGMK holds promise for protein engineering because it contains a combination of hydrophobic and hydrophilic residues. Designing a protein with this sequence can be approached using computational techniques to forecast its secondary and tertiary structures, though challenges arise from the intricacies of protein folding and the computational resources required. It is essential to validate computational predictions to confirm their reliability and applicability in real-world scenarios. Methods like cross-validation, regularization, early stopping, ensemble learning, and transfer learning are valuable for avoiding overfitting and underfitting. Additionally, interpreting and elucidating model results can enhance practical implementation.}
\end{quote}

We can further delve into detailed steps for sequence engineering by prompting the model to offer more targeted design recommendations aimed at improving protein stability:

\begin{quote}
\small\texttt{What modifications could be made to the following protein sequence to enhance its stability: MSKSNTYRMLVLLEDDTKINKEDEKFLKGKPGKMHEFVDELILPFNVDELDELNTWFDKFDAEICIPYVEILKESGMK?}
\end{quote}

A summary of the discussion (full transcript available in Supplementary Material, SI Conv Example 2):

\begin{quote}
\small\texttt{The dialogue focuses on strategies to increase protein stability by incorporating hydrophobic amino acids and disulfide bonds. The biologist explains that hydrophobic residues generally contribute to stability and that disulfide bonds can also strengthen the protein structure. The engineer requests specific examples of hydrophobic amino acids and advice on the appropriate number of disulfide bonds to include. The biologist cites proteins such as insulin, bovine pancreatic trypsin inhibitor, and green fluorescent protein as successful cases where these modifications have been employed. Potential drawbacks are also addressed, including possible alterations in protein function, increased structural complexity, and risks of aggregation. The biologist recommends overcoming these challenges through judicious amino acid selection, computational modeling, and empirical validation.}
\end{quote}

The summary of the main points:

\begin{quote}
\small\texttt{\begin{itemize}
\item To enhance the stability of a protein sequence, one can incorporate hydrophobic amino acids and disulfide bonds.

Hydrophobic amino acids generally increase protein stability, while disulfide bonds also contribute to enhanced stability.

Proteins that have been effectively stabilized using these modifications include insulin, bovine pancreatic trypsin inhibitor, and green fluorescent protein.
\item Possible drawbacks of adding hydrophobic amino acids or disulfide bonds include alterations in protein function, increased complexity, and a risk of aggregation.
\item These challenges can be mitigated through careful amino acid selection, computational modeling, and experimental verification.
\end{itemize}
}
\end{quote}

The main conclusion is:

\begin{quote}
\small\texttt{The key point from the discussion is that to increase the stability of the protein sequence MSKSNTYRMLVLLEDDTKINKEDEKFLKGKPGKMHEFVDELILPFNVDELDELNTWFDKFDAEICIPYVEILKESGMK, introducing hydrophobic amino acids and disulfide bonds is effective. This directly addresses the initial question.}
\end{quote}

From the viewpoint of protein engineering, these recommendations are logical as they are likely to yield more stable protein designs. Specific strategies and examples are provided that researchers can further investigate.

In a different instance, three sequences are considered. The model is first asked to evaluate the stability of each, then to analyze the results: identify the strongest candidate and explain why that protein is probably the most stable. Additionally, the model is asked to describe the experimental steps required to produce it in the laboratory. The three sequences (manually generated by modifying one from the test set) are:

\begin{quote}
\small\texttt{
\begin{itemize}
    \item LNTWFDKFDAEICIPYVEILKESGMK \item AITWFDKFDAEICIPYVEALKIAAMK \item AAAAAAAAKFDAAEICIIAAAAAAAAKIAAMK
\end{itemize}
}
\end{quote}

This example evaluates whether the model can integrate multiple skills and adapt dynamically as the conversation advances. That is, to flexibly switch between protein mechanics, logical reasoning, hypothesis generation, and similar capabilities.

The dialogue is summarized in Figure~\ref{fig:conversation_evolve}, accompanied by an illustration of how the scalings evolve throughout the conversation (the analysis includes both the detailed layer-wise adapter scaling and an aggregated analysis to provide a measure of the effective utilization of the different experts). The findings clearly demonstrate that the protein task is initially highly pertinent, gradually transitioning to a stronger representation of the bioinspired adapter and the biology adapter. By the end of the conversation, the biology adapter becomes the most dominant.

To evaluate the predictions and assess the quality of the responses, the three proteins considered in this task are shown in Figure~\ref{fig:protein_visuals} (proteins folded using Alpha Fold 2~\cite{Jumper2021HighlyAlphaFoldb}, via ColabFold \cite{Mirdita2022ColabFold:All}). We observe that the most stable structure in the center exhibits the most organized secondary structure, consistent with the idea that it produces the highest unfolding force, as predicted by the X-LoRA model in Figure~\ref{fig:conversation_evolve}. The other two structures are less organized; the first is predominantly helical with some turns, and the last is partly unstructured and has a kink in its geometry. These characteristics account for the lowest predicted unfolding force according to the model.

By examining the visual depictions of the folded structures, we verify that the most stable structure located in the center exhibits the most orderly secondary organization, consistent with the idea that it produces the highest unfolding force as forecasted by the model. We observe that the other two models are less structured. The first model is predominantly helical with some turns, while the last model is partly unstructured and displays a kink in its geometry. We suggest that these characteristics likely account for the lowest observed unfolding force predicted by the model, aligning with the X-LoRA model’s predictions.

In a subsequent example, we concentrate on the energetic characteristics of proteins, specifically analyzing the unfolding energy as the protein is stretched from its initial configuration to a fully unfolded conformation~\cite{Ni2023ForceGen:Model}. We once again consider three sequences (the first two are the same as those examined previously, while the third is a longer sequence composed of a combination of the second sequence at the beginning and end and the first sequence in the center):

\begin{quote}
\small\texttt{\begin{itemize}
    \item LNTWFDKFDAEICIPYVEILKESGMK \item AITWFDKFDAEICIPYVEALKIAAMK \item AITWFDKFDAEICIPYVEALKIAAMKLNTWFDKFDAEICIPYVEILKESGMKAITWFDKFDAEICIPYVEALKIAAMK
\end{itemize}
}
\end{quote}

We initially instruct the model to determine the unfolding energy of each sequence and subsequently to analyze the results: Based on the amino acid sequences, we request an explanation as to why this protein is probably the most stable. Finally, we ask for possible functions that the most stable protein might perform.

Figure~\ref{fig:MJB_20} presents a transcript of a dialogue that combines different expert knowledge and synthesizes information across these domains, demonstrated here through the analysis of the unfolding energy of three proteins. By the conclusion of the dialogue, the CoT/reasoning adapter emerges as the dominant component, with the X-LoRA model addressing a query that involves reasoning over the results to forecast probable protein structural characteristics as well as potential protein functions. Specifically, the mechanistic inquiry regarding the source of protein stability highlights the role of hydrophobic interactions among nonpolar amino acids and hydrogen bonding among polar amino acids. The model predicts that multiple hydrophobic amino acids (I, F, W, Y) and hydrogen-bonding amino acids (D, E, K) contribute to the protein’s stability.

Additionally, the model forecasts that cysteine (C) residues in the sequence will promote the formation of disulfide bonds, which further enhance the stability of the protein’s structure. Notably, the critical predictions concerning structural features that explain the high stability are accurately identified, as confirmed by structural analyses displayed in Fig.~\ref{fig:MJB_21}. This analysis validates the presence of all these features. It also clarifies why the other two proteins, which are shorter segments exhibiting variable helical domain stability, have substantially lower unfolding energies.

\subsection{Adversarial agentic modeling to connect distinct scholarly disciplines and knowledge yielding ontological knowledge graph generation}

Given that our X-LoRA model possesses extensive knowledge across multiple fields, we can leverage it to explore links among diverse concepts, knowledge repositories, and specialties. In one experiment, we posed two queries to the model, each structured similarly but emphasizing different aspects. Subsequently, we generated triplets to construct an ontological knowledge graph that organizes the answers into more structured representations~\cite{Giesa2012CategoryDesign,Spivak2011ReoccurringAnalogies}. The resulting graph offers an integrated perspective on the insights produced and visually depicts relationships between concepts in a clear and mechanistic way.

We employ two opposing agents to address the question. Agent 1 is assigned the role of a "critical physicist," while Agent 2 functions as a "philosopher." Agent 1 initiates the inquiry and is responsible for probing further and generating follow-up questions based on each of Agent 2's replies. As the dialogue progresses, it delves more deeply into the topic, revealing multiple facets. The chosen question is deliberately designed to span several domains of disciplinary knowledge, aiming to test the model's capacity to synthesize ideas and concepts.

The question posed is:

\begin{quote}
\small\texttt{Consider an abstract representation of a classical musical composition as a graph of interacting notes, chords, melodies, rhythms and so on.

What happens if we would apply pressure, up to a point of failure, on this abstract concept of music?

Provide logical steps, mathematical reasoning, and discuss specifically the physics of the problem.}
\end{quote}

The complete conversation is included in the Supplementary Material (refer to SI Conv Example 3). For brevity, we present only the summary and key insights here, which encapsulate the main points:

\begin{quote}
\small\texttt{The discussion centers on the effect of applying pressure to an abstract depiction of a classical musical composition, viewed as a physical system comprising interacting notes, chords, melodies, and rhythms. Pressure is quantified through stress, with the failure threshold identified when stress reaches a critical level. To validate the mathematical model's predictions, a physical analog of the abstract musical composition can be constructed, and experimental outcomes compared against theoretical expectations. Limitations of the proposed mathematical model include simplified interaction representations, static nature, uniform distribution assumptions, and linear stress-pressure relationships. To enhance model accuracy and validity, incorporation of more detailed and realistic models, dynamic considerations, heterogeneous distributions, and nonlinear terms is recommended.

[...]

The primary takeaway is that applying pressure to an abstract representation of a classical musical composition can be conceptualized as modifying the force field among interacting notes, chords, melodies, and rhythms. With increased pressure, system stress rises, and at a critical threshold, the system fails and collapses under its own weight. This understanding addresses the original question by offering a physical rationale for how pressure influences the abstract notion of music and how this influence can be quantified via stress.}
\end{quote}

Although the generated text yields valuable insights, constructing an ontological representation of the data can aid in interpretable analysis~\cite{Buehler2023MechGPTModalities}. For graph creation, we combine the entire dialogue and build a knowledge graph. An example of such a graph is depicted in Figure~\ref{fig:graph}. The resultant graph is partitioned into communities using the Girvan-Newman algorithm~\cite{Girvan2002CommunityNetworks}, identifying a total of 12 communities. The figure displays both the overall graph (Figure~\ref{fig:graph}a) and more detailed views with node labels (Figure~\ref{fig:graph}b-c).

\subsection{Creation of X-LoRA-Gemma Integrating Protein, Chemical, Bio-inspired, and Materials Mechanics Features}
\label{gemma-section}

To demonstrate that the proposed method is effective with other foundational models possessing different architectures, we trained an alternative X-LoRA model, this time utilizing the Gemma-7B-it model~\cite{Google/gemma-7b-itFace}. This version incorporates four LoRA adapters, defined as follows:

\begin{enumerate}
    \item Bioinspired materials \item Mechanics and materials \item Protein mechanics tasks (including generative sequence-to-property and inverse design capabilities) \item Quantum-mechanics based molecular properties QM9 (with generative SMILES-to-property and inverse functionalities; see Table~\ref{tab:table_QM9} for definitions and an overview of all properties calculated in this dataset)~\cite{Ruddigkeit2012EnumerationGDB-17,Ramakrishnan2015ElectronicSpace}
\end{enumerate}

The X-LoRA-Gemma model is developed similarly to the initial X-LoRA model but is designed to incorporate diverse capabilities, notably the ability to predict a set of 12 quantum mechanical properties and to generate molecules tailored to these 12 quantum mechanical property targets. This new adapter is trained alongside the other adapters mentioned above, which use the same datasets as those employed in the previous experiments. It is important to note that while the majority of the paper concentrates on the Zephyr-based X-LoRA model, in this section we differentiate between the two models by referring to the original as X-LoRA-Zephyr and the Gemma-based variant as X-LoRA-Gemma.

Figure~\ref{fig:protein_force_energy_prediction} presents the overall performance of X-LoRA, highlighting predictions of unfolding force and unfolding energy from sequence data ($R_2 = 0.87$), as well as predictions of quantum mechanical properties of molecules from the QM9 dataset ($R_2 = 0.96$). Both tasks show excellent results, with the protein design task performance being slightly superior to that of the original X-LoRA model.

Next, we illustrate the functionality of X-LoRA-Gemma through a straightforward multi-agent design experiment. We focus on a small organic molecule, \texttt{CCCC12CC1C(=O)N2} (extracted from the QM9 dataset test set). This bicyclic molecule contains a ketone (=O) and an amine (N) within the ring structure (see Figure~\ref{fig:QM9_design}(a) for a visualization). We pose the following question to the X-LoRA-Gemma model:

\begin{quote}
\small\texttt{Here is a molecule in SMILES code: CCCC12CC1C(=O)N2

What modifications would improve this molecule’s suitability for polymer design? }
\end{quote}

The response is:

\begin{quote}
\small\texttt{To enhance the usefulness of this molecule for polymer synthesis, you might:
\begin{enumerate}
\item  Add a reactive functional group, like an amine or hydroxyl group, to facilitate further modification and crosslinking.
\item  Extend the carbon backbone to be longer and more flexible, improving the molecule's capacity to form polymers with varied structures and characteristics.
\item  Attach a side chain containing a functional group, such as a carboxyl group, enabling the formation of covalent bonds with other molecules and polymers.
\item  Integrate a cyclic component into the molecule, for example, a lactam ring, to produce a polymer that is more rigid and stable.
\item  Incorporate diverse side chains to generate a polymer with a broad range of properties and functionalities.
\item  Vary the number of carbons in the backbone to permit the synthesis of polymers with different molecular weights and attributes.
\item  Introduce a crosslinking site within the molecule to yield a polymer that is more stable and less reactive.
\item  Design a repeating unit structure to achieve a polymer with defined molecular weight and characteristics.
\item  Combine several of these approaches to develop a flexible and functional polymer suited for multiple applications.
\end{enumerate}}
\end{quote}

Next, we concentrate on the first recommendation.

\begin{quote}
\small\texttt{What approaches can I take to enhance the overall reactivity of a molecule? I intend to utilize this molecule to construct a polymer material.

The molecule under consideration, represented in SMILES notation, is: CCCC12CC1C(=O)N2

I am able to adjust the following parameters:

Dipole moment, mu Polarizability, alpha Highest Occupied Molecular Orbital energy, homo Lowest Unoccupied Molecular Orbital energy, lumo HOMO-LUMO gap: The energy difference between HOMO and LUMO, gap Electronic spatial extent Zero-point vibrational energy Heat capacity at constant volume Internal energy at 0 K Internal energy at 298.15 K Enthalpy at 298.15 K Free energy at 298.15 K Which are the top THREE properties I should modify, and should I increase or decrease them?}
\end{quote}

The feedback from X-LoRA-Gemma is as follows:

\begin{quote}
\small\texttt{\begin{itemize}
    \item  HOMO-LUMO gap (reduction) \item  Dipole moment (enhancement) \item  Polarizability (enhancement)
\end{itemize}
Adjusting these properties can boost the overall reactivity of the molecule, making it better suited for applications in polymer material construction. }
\end{quote}

We leverage both the generative and predictive functionalities of X-LoRA-Gemma to engineer these specific molecular property modifications.

For the optimization, we employ a two-agent framework composed of two separate X-LoRA-Gemma instances. One agent performs the generative function \texttt{GenerateMolecularProperties<...>} to propose candidate molecules. The other agent carries out the evaluation via \texttt{CalculateMolecularProperties<...>} and quantifies the error between the desired and predicted properties. This cycle continues iteratively until either the maximum number of iterations is reached or the mean-squared error falls below a predefined threshold.

Figure~\ref{fig:QM9_design} presents the outcomes of this entire workflow. Panel (a) illustrates the target design and its deviation relative to the initial molecule’s properties. Panel (b) displays the top 20 candidate molecules generated by the algorithm, accompanied by the mean-squared error distribution (c) and a plot showing the mean-squared error progression over the generated candidates (d). The optimal design identified is \texttt{CC(C)C1=CC(=O)NC=C1}, which closely matches the targeted properties as depicted in panel (e). Finally, panel (f) provides an optimized 3D molecular structure, obtained via UFF~\cite{Rappe1992UFFSimulations} force field optimization.

To evaluate the molecule’s reactivity, we also calculated the Gasteiger charges~\cite{Gasteiger1980IterativeCharges} (see Figure~\ref{fig:QM9_design}(f) for a visual summary). The molecule contains an aromatic ring with a ketone group (C=O), which is more electrophilic than typical C–C or C–H bonds, rendering it prone to nucleophilic attack. The combination of the nitrogen atom in the aromatic ring and the ketone group forms a pyridone ring system. This ring exhibits higher reactivity than a simple aromatic ring due to the heteroatom (nitrogen) and the carbonyl group, which provide reactive sites such as the carbonyl carbon for nucleophilic addition and positions on the ring for electrophilic substitution. The nitrogen, as part of the heteroaromatic system, influences both the aromaticity and the electron density distribution throughout the ring, thereby modulating its chemical reactivity. This modification makes the nitrogen-containing ring a reactive site, particularly susceptible to electrophilic attacks driven by the electron-donating effect of the nitrogen atom.

The carbon atom within the ring, which is positioned next to both an oxygen and a nitrogen atom and carries a partial positive charge (+0.25), indicates a degree of electron deficiency. This occurs because oxygen, being more electronegative, draws electron density toward itself, thereby reducing the electron density on the carbon. This effect is somewhat moderated by the neighboring nitrogen, which has an electronegativity lower than oxygen’s but higher than carbon’s. Nonetheless, the partial positive charge on the carbon signifies it as a probable site for nucleophilic attack, where nucleophiles (electron-rich species) are attracted to this electron-poor carbon center.

The nitrogen atom, with a partial negative charge of -0.33, reflects that it is relatively electron-rich compared to its standard state. This enrichment can be attributed to its lone pair of electrons and the influence of the aromatic ring’s electron system. In aromatic heterocycles, the lone pair on nitrogen may participate in the delocalized $\pi$-electron system; however, the degree of this involvement depends on the ring’s configuration and the electronegativity of nearby atoms. The partial negative charge suggests that nitrogen is more nucleophilic, implying a greater tendency to donate electrons, either by forming bonds with electrophiles or undergoing protonation. Additional quantum mechanical studies should be conducted to explore this behavior further.

The carbon bearing a partial positive charge represents a reactive site susceptible to nucleophilic attacks. Various molecules could interact at this position through mechanisms involving nucleophiles targeting carbons deficient in electron density.

The hydrogen atom attached to the nitrogen, exhibiting a partial positive charge of +0.17, is a prime candidate for engaging in hydrogen bonding with other molecules or atoms. Hydrogen bonds are a form of dipole-dipole interaction that arises when a hydrogen atom bonded to a highly electronegative atom (such as nitrogen, oxygen, or fluorine) interacts with another electronegative atom possessing a lone electron pair. This characteristic endows the molecule with the ability to act as a hydrogen bond donor, an important factor influencing the molecule’s solubility in water and other polar solvents, as well as its interactions within biological environments.

The ketone group positioned next to the nitrogen-containing aromatic ring provides distinctive reactivity that can be harnessed in polymer synthesis~\cite{McQuarrieSimonPhysicalChemistry,CareySundbergAdvancedOrganicChemistry,Varghese2022BeyondApplications,Varghese2022BeyondApplications}. This compound is classified as a lactam, due to the cyclic amide formed by the ketone and nitrogen within the ring. Lactams play a significant role in polymer chemistry, especially in the creation of polyamides, a group of polymers widely utilized in material applications. The ketone may participate in several chemical reactions relevant to polymer chemistry, such as nucleophilic addition, which can be exploited to lengthen the polymer chain or to add side chains that alter the polymer’s characteristics. The nitrogen atom in the aromatic ring can notably influence the polymer’s electronic attributes and engage in hydrogen bonding, which may affect properties like melting point, solubility, and mechanical strength. N-heteroaromatic compounds also exhibit the ability to coordinate with metal ions, potentially enabling uses in catalysis or in materials with electronic or photonic functions. Additionally, this molecule could undergo polymerization via its amide (lactam) group; polyamides are recognized for their durability, elasticity, and chemical resistance, making them suitable for diverse applications ranging from textiles (such as nylon) to automotive components~\cite{Varghese2022BeyondApplications,Varghese2022BeyondApplications}.

Alkyl substituents can enhance the hydrophobic nature of the molecule, so we expect that polymers synthesized from monomers containing such alkyl groups will exhibit increased hydrophobicity. This alteration can influence their solubility in various solvents as well as their interactions with water. Such characteristics may be beneficial for applications that demand water resistance or particular solubility profiles. Additionally, alkyl groups can impact the mechanical behavior of the polymer. For example, they may serve as plasticizers within the polymer matrix, thereby enhancing the polymer's flexibility. This effect could result in materials that are more flexible or possess improved impact resistance, depending on the alkyl groups' length and branching.

\subsubsection{Additional benchmarking} We present further benchmarking results for the X-LoRA-Gemma model, illustrated in Figure~\ref{fig:perf_MMLU}, based on a subset of the MMLU benchmark that targets chemistry, physics, and biology domains (which closely align with the areas on which our model was trained)~\cite{Hendrycks2020MeasuringUnderstanding}. Our findings show that both X-LoRA models evaluated in this study outperform their respective base models. Specifically, the X-LoRA models achieve higher scores than the base versions (Zephyr model: 54.6\% for the base versus 56.6\% for the X-LoRA model; Gemma model: 40.6\% for the base and 52.4\% for the X-LoRA-Gemma model). The X-LoRA-Zephyr model delivers the highest overall performance, although both models perform comparably well.

\section{Conclusion}

Multimodal LLMs have numerous current and prospective applications in scientific fields, accompanied by a significant demand for specialized, domain-specific expertise. Beyond general language processing tasks, the capability to perform computational or generative tasks, as illustrated here with protein mechanics, is essential~\cite{Luu2023BioinspiredLLM:Materials,Buehler2023MechGPTModalities,Buehler2023MeLMProblemsb}. X-LoRA, which is built on open-source language models, allows for the flexible integration of expertise and knowledge across various domains, as demonstrated by multiple examples, including Figure~\ref{fig:conversation_evolve}. This investigation revealed how the X-LoRA model autonomously and adaptively adjusts its use of adapters, blending certain adapters in a sophisticated manner throughout the LoRA layers. Such intricate mixing of LoRA layers is consistently observed across all examples examined here, such as in Figures~\ref{fig:QA_weights_sample_0} and \ref{fig:QA_weights}. This also presents an intriguing challenge to better comprehend the fundamental mechanisms from a physics standpoint, where the LLM model can be conceptualized as a large-scale graph-forming system that leverages complex graph structures. Consequently, the way adaptation experts are mixed at deep layer scales, as performed here, modifies these mechanisms to facilitate the development of appropriate solutions.

The introduced algorithm enhances the traditional inference method, which uses a single forward pass, by incorporating two forward passes. This trains the model to initially `contemplate' the question and how it might reconfigure itself before providing an answer. This approach implements a straightforward form of `self-awareness,’ enabling the model to modify its own architecture to optimally address a task. Consequently, despite its relatively modest size of 7 billion parameters, the model can reason across a wide range of scientific disciplines (including biological materials, mathematics, physics, chemistry, logic, mechanics, etc.), greatly boosting its ability to generate novel solutions. Furthermore, this may inspire the design of alternative model architectures, such as methods that extend the reconfigurational strategy to portions of the model or to the entire model in non-adapted systems. In such cases, the scaling head would dynamically reconfigure parts of a trained model or integrate multiple models. As demonstrated by X-LoRA, these strategies can serve as powerful tools to expand or combine the capabilities of various models.

We conducted multiple performance evaluations, including a comparison with significantly larger models (Figure~\ref{fig:perf_comp_bioinspiredexam}), where X-LoRA clearly demonstrated superior performance. The detailed comparison presented in Table~\ref{tab:table_qa} revealed that X-LoRA significantly outperforms the base model across a wide range of tasks and application areas. Additional comparisons involved quantitative evaluations of numerical prediction tasks, such as protein mechanics and quantum-mechanical molecular properties, where X-LoRA achieved high accuracy with $R_2$ values ranging from 0.85 to 0.96. For reference, these results surpassed previously reported outcomes, such as predicting protein unfolding force using the MeLM model~\cite{Buehler2023MeLMProblemsb}, which yielded $R_2=0.78$. The ability to accurately solve multiple prediction tasks, coupled with inverse design tasks and complex reasoning capabilities, was highlighted in the molecular design example (Section~\ref{gemma-section}). This case also demonstrated that X-LoRA can be effectively implemented using different base model architectures (Gemma versus Zephyr/Mistral). Future research could extend this approach and develop X-LoRA models for larger base models.

A limitation of the proposed method is the need for two forward passes, potentially increasing computational costs: one to compute hidden states serving as inputs to the X-LoRA scaling head, and another with $\Lambda$ applied to the adapters. However, the first pass is leveraged by the X-LoRA scaling head and thus utilizes fundamental LLM capabilities, rather than requiring separate knowledge acquisition through training a larger scaling head. For efficient deployment with model-serving tools, separate key-value caches need to be maintained for both the scaling and forward passes. To enhance inference speed, we developed \texttt{Mistral.rs}, an LLM serving platform implemented in Rust~\cite{matsakis2014rust} that incorporates X-LoRA and several performance optimization techniques. Nonetheless, the main benefit of our approach is its compatibility with any model without needing to restructure internal logic. We consider this a significant advantage, especially since it enables leveraging the extensive resources available in the Hugging Face ecosystem.

A key benefit of the proposed method is its straightforward implementation within any existing LLM (for example, with the release of the X-LoRA code, it can readily be applied to any model in the Hugging Face ecosystem). Another benefit is that the scaling head leverages the inherent capabilities of the underlying LLM, learning nuanced strategies for optimally combining adapters to produce certain answers. During X-LoRA training, it is ideal to use complex question-answer pairs or dialogues that enable the model to discover the most effective combinatorial approach. We conducted multiple analyses related to protein design and mechanics. For example, Figure~\ref{fig:MJB_20} and Fig.~\ref{fig:MJB_21} show a stability assessment of various protein sequences, including a comparative analysis, reasoning about predicted outcomes, and a mechanistic explanation offering a fundamental chemical basis for the observed behavior. As demonstrated in Fig.~\ref{fig:MJB_21}, this was corroborated by a structural examination of the protein’s folded 3D conformation.

Regarding future directions, further studies could investigate a broader range of adapter experts, extending beyond the initial expertise present in our LoRA collection. The X-LoRA model surpasses the performance of either the base model alone or the base model combined with a single adapter, delivering answers to intricate questions by utilizing specialized knowledge contributed by the adapter set. There also appear to be interactions among the various adapters, as suggested by the complex mixing patterns of scaling weights across layers. This phenomenon warrants additional investigation and comparison with techniques such as SLERP, which merge models via different approaches. While we trained the X-LoRA scaling head using a subset of training data—comprising a few hundred samples drawn from the original datasets that informed the adapters’ development—a more targeted creation of training sets tailored for this task could be formulated. For example, we observed that multiple-choice questions tend to engage reasoning-focused adapters, reflecting their training on similar datasets. Extra prompting is often necessary to guide the scaling head toward understanding the relevant domain for particular queries, potentially through specific system messages or other mechanisms. Broadly, the design of suitable training datasets is a promising area for further research, including methodologies that encourage effective mixing through layer-wise scaling. Our experiments tracking these scaling dynamics over extended dialogues indicate that the model can adaptively adjust scaling strategies to best address particular tasks, highlighting a promising capability.

We identified intriguing patterns regarding the activation of experts across layers, with valuable insights explored in the paper. This analysis could be further developed and potentially provide significant understanding of how and why mixing components of models or groups of layers might be beneficial. Such exploration would be particularly compelling when applying the X-LoRA technique beyond protein mechanics and protein design, as demonstrated here. We reserve this for future research.

Another critical direction is the continued investigation of agentic modeling, as exemplified here through the use of two X-LoRA models engaging in adversarial interaction. Subsequent work could involve more than two models to enrich interactions and introduce additional functionalities, driving models toward novel generative domains. This approach may also facilitate the creation of methods that incorporate physics-based modeling or other validation steps, leveraging code-writing/executing agents or agents utilizing function calling or other processing strategies~\cite{Ghafarollahi2024ProtAgents:Learning}.

\section{Materials and Methods}

The X-LoRA code and example implementations are accessible at \url{https://github.com/EricLBuehler/xlora}. Model weights and datasets associated with the X-LoRA referenced in this paper, along with further examples, can be found at \url{https://huggingface.co/lamm-mit/x-lora}. The X-LoRA-Gemma model is hosted at \url{https://huggingface.co/lamm-mit/x-lora-gemma-7b}. The \texttt{Mistral.rs} inference engine is available at \url{https://github.com/EricLBuehler/mistral.rs}.

\subsection{Mathematical formulation of X-LoRA} We define the $\mathbf{scalings}$ as a matrix containing coefficients at both token and layer levels for LoRA adapters. The X-LoRA $\mathbf{scaling}$ $\mathbf{head}$ processes hidden states from the base model through linear fully-connected layers to generate scalings, followed by a softmax operation on the final dimension to guarantee that scalings across all experts sum to one. ReLU activation functions and dropout layers are interspersed between the linear fully-connected layers within the X-LoRA scaling head. The size of the hidden layers can be adjusted, either to effectively reduce the dimension from the hidden space to the scaling dimension or to implement a widening and narrowing structure analogous to the feed-forward layer in transformer blocks.

The scalings constitute a matrix $\Lambda \in \mathbb{R}^{s \times l \times n}$, where $s$, $l$, and $n$ denote the sequence length, number of layers, and number of adapters, respectively. For each layer, the relevant scalings are extracted as a matrix $\lambda \in \mathbb{R}^{s \times n}$. These scalings $\lambda$ are applied to each LoRA adapter by broadcasting and multiplying the corresponding X-LoRA scaling $\lambda_i \in \mathbb{R}^{1 \times 1 \times s}$ to the input of the decomposition matrices $A$ and $B$.

Given $W_0 \in \mathbb{R}^{d \times k}$, $B_i \in \mathbb{R}^{d \times r_i}$, and $A_i \in \mathbb{R}^{r_i \times d}$ where the rank $r_i \ll \min(d,k)$ is specific to each adapter $i$, the forward pass of an X-LoRA layer is formulated by:

\begin{equation}
    h = W_0 x + \sum_{i=1}^{n} B_i A_i (x \cdot \lambda_i) \alpha_i
\end{equation}

Representations of the scalings are shown in the main paper, visualized either through bar plots or by using the \texttt{contourf} function from Matplotlib~\cite{MatplotlibDocumentation}.

\subsection{X-LoRA implementation algorithm and code} The selected implementation prioritizes user-friendliness and the flexibility to adapt the method to a wide range of models, particularly those in the Hugging Face ecosystem. The X-LoRA algorithm functions by:

\begin{enumerate}
    \item Keeping the base model and adapters frozen while tuning performance.
    \item Combining LoRA adapters based on the predictions from the X-LoRA scaling head.
\end{enumerate}

All code is developed in Python and PyTorch~\cite{Paszke2019PyTorchLibrary}.

\subsubsection{Dual forward pass strategy for self-aware inference} Since X-LoRA needs to compute hidden states to predict scaling factors, we employ a dual forward pass technique. The first pass enables the model to analyze the task and determine how to adjust itself, followed by the actual inference step. This approach introduces a form of `self-awareness' that allows X-LoRA to allocate additional computation for reflection instead of immediate response.

We refer to the first pass as the scaling pass and the second pass as the forward pass. During the scaling pass, we initialize the scalings $\Lambda$ to zero. We also tested initializing $\Lambda$ as ${n}^{-1}$ and observed comparable outcomes, demonstrating stable performance (in the X-LoRA package, this setting can be adjusted via configuration). To facilitate straightforward application of X-LoRA across various models, we incorporate a forward pass hook. This mechanism is crucial to coordinating the forward and scaling passes and constitutes the primary feature enabling compatibility with all models.

\begin{enumerate}
    \item $\mathbf{Scaling}$ $\mathbf{pass}$: Conducted by running the base model with adapters fixed at constant values. The hidden states produced are then used by the X-LoRA scaling head to compute $\Lambda$.
    \item $\mathbf{Forward}$ $\mathbf{pass}$: Uses $\Lambda$ to blend the adapters during the model's forward pass. The output logits serve as the final output of the X-LoRA model.
\end{enumerate}

It is important to note that sharing the same key-value cache between the scaling and forward passes may disrupt the correct input flow to the scaling head. This issue arises because the key-value cache would persist and accumulate data across both forward passes. Therefore, we disable the key-value cache during training and inference and have implemented appropriate management of this in \texttt{Mistral.rs}.

\subsubsection{Freezing weights and selectively enabling model segments for training} The base model and adapters remain frozen (meaning their weights are not updated during training), so the only trainable component in an X-LoRA model is the X-LoRA scaling head. However, the provided code allows for the possibility to unfreeze all adapters alongside the X-LoRA scaling head. It is also possible to train the entire model simultaneously—that is, the X-LoRA scaling head, adapters, and the foundational base model—either with uniform or distinct learning rates. Furthermore, future developments could extend these ideas to create scaling functions that operate between two or more foundation models, facilitating a transparent implementation of SLERP-like modeling. This direction could be pursued in subsequent research.

\subsubsection{X-LoRA layers} An X-LoRA layer functions by scaling and then combining outputs from multiple adapters. Each X-LoRA layer encapsulates the original LoRA layer and thus can access its corresponding adapters.

To apply X-LoRA effectively to any model, the algorithm replaces the original LoRA layers with X-LoRA layers and injects the new forward pass mechanics into these layers. This approach exploits features of Python's object system to allow seamless substitution.

In brief, during the forward pass, each X-LoRA layer performs the following operations (refer to Figure \ref{fig:general_arch}):

\begin{enumerate}
    \item $\mathbf{Scaling}$: Scale the input signal of each LoRA adapter by the factor $\lambda_i$.
    \item $\mathbf{Forward}$: Execute the forward pass of the encapsulated LoRA layer.
\end{enumerate}

\subsection{Efficient Inference with \texttt{Mistral.rs} implemented in Rust} To accelerate inference, we created \texttt{Mistral.rs}, a LLM serving platform written in Rust~\cite{matsakis2014rust}, which incorporates X-LoRA and various performance optimizations. Specifically, the following techniques have been integrated, presented here in no particular order:

\begin{itemize}
    \item LoRA adapter weight stacking: When all $A_i$ and $B_i$ matrices across LoRA adapters share identical dimensions, their weight matrices $A_i$ and $B_i$ can be stacked along the first dimension.

Mathematically, given $W_0 \in \mathbb{R}^{d \times k}$, for adapter matrices $A_i \in \mathbb{R}^{r \times d}$ and $B_i \in \mathbb{R}^{d \times r}$ where $r_i \ll {\rm min} (d,k)$, we stack the matrices so that $A \in \mathbb{R}^{i \times r \times d}$ and $B \in \mathbb{R}^{i \times d \times r}$. The parameters $\alpha_i$ are also applied to the concatenated matrices during initialization.

To apply scalings, we rearrange the dimensions so that $\lambda \in \mathbb{R}^{i \times 1 \times 1}$. Afterwards, these are applied using broadcast multiplication.

\item Non-granular scalings To minimize the number of times the scaling pass is executed, we provide an option for non-granular scalings. This approach caches the scalings after $n$ completed runs, leveraging the KV cache to maintain the sequence length at one for the rest of the request. This method significantly lowers latency by reducing how often the base model is invoked.

\item Fused Compute Unified Device Architecture (CUDA) kernels For optimizing the forward pass, we implement fused CUDA kernels. Such fused kernels allow complex operations, like Rotary Embedding (\cite{su2021roformer}), which would normally be performed sequentially, to be merged into a single kernel and executed in parallel on a GPU. Our kernels are based on the vLLM implementation (\cite{kwon2023efficient}).

\item Quantization The inference engine supports quantization, enabling the use of a quantized base model. Initial experiments indicate that quantization down to 8 bits performs well, even for models originally trained with \texttt{bfloat16}, for example.

\end{itemize}

\subsection{Datasets}

Each adapter is trained using specialized datasets, summarized in Table~\ref{tab:table_loradata}, which includes references to the original sources and literature.

\subsection{Training strategy} Training proceeds in phases. We utilize the Zephyr-7B-$\beta$ model~\cite{HuggingFaceH4/zephyr-7b-betaFace}, which is built upon the Mistral-7B model\cite{Jiang2023Mistral7B}, and train multiple adapters with the datasets outlined previously. The Zephyr-7B-$\beta$ tokenizer is employed.

The first eight adapters are trained with question-answer pairs as provided by the respective datasets (see Table~\ref{tab:table_loradata}).

Protein mechanics tasks are trained in a similar manner but employ specific forward and inverse instruction sets. The bidirectional instruction tasks include:

\begin{quote}
\tiny {\texttt{CalculateForce< G E E C D C G S P ....> [0.310]\\
CalculateEnergy< G E E C D C G S P ....> [0.220]\\ 
CalculateForceHistory< G E E C D C G S P ....> [0.004,0.034,0.125,0.142,0.159,0.102,0.079,0.073,0.131, ...]\\
GenerateForce<0.310>\,[ G E E C D C G S P ....]\\
GenerateEnergy<0.220>\,[ G E E C D C G S P ....]\\
GenerateForceHistory<0.004,0.034,0.125,0.142,0.159,0.102,0.079,0.073,0.131, ...>\,[ G E E C D C G S P ....] }}
\end{quote}

These denote three distinct tasks: calculating the peak unfolding force of a protein, determining the unfolding energy, and obtaining the force-deformation history. Each of these forward tasks is paired with an inverse task of the same kind, which instead outputs a protein sequence, resulting in a total of six tasks.

Each LoRA adapter has a rank set to $r=16$ with $\alpha=16$. The training of each adapter involves five target modules: \texttt{v\_proj}, \texttt{k\_proj}, \texttt{q\_proj}, \texttt{gate\_proj}, and \texttt{o\_proj}, applying a dropout rate of 0.05. Typically, LoRA adapters are trained using relatively low ranks, ranging from 4 to 32, since increasing the rank generally does not enhance performance but does raise computational costs; a lower rank yields a more compact model with fewer parameters~\cite{hu2021lora}. The selected values for $r$, $\alpha$, and other parameters follow commonly adopted settings in the literature that have demonstrated effective results. Additionally, previous research conducted by the author has tested different configurations and concluded that the parameters chosen here are a suitable compromise~\cite{Luu2023BioinspiredLLM:Materials}.

The X-LoRA scaling head is trained using approximately 2,000 samples, which are a subset taken from the original dataset employed for training the adapters (we also augmented the training data with additional GPT-4 generated multiple choice question samples). The optimizer used is \texttt{paged\_adamw\_8bit} with gradient clipping set to 0.3, a learning rate of $2\times 10^{-4}$ with warmup, and four gradient accumulation steps, using a batch size of one. Training of the X-LoRA model is conducted for roughly 10,000 steps. The loss at the end of training is approximately 0.28.

Given that the underlying Mistral architecture used here as the base model has 32 layers, and we incorporate adapters into 5 of the transformer layers each, this results in a total of 160 LoRA layers, each of which is replaced by an X-LoRA layer. The X-LoRA scaling head generates scaling factors for each of these 160 LoRA layers (with nine LoRA experts, this yields a total of $160\times 9=1440$ $\lambda_i$ values computed per token). The model employs a single linear layer with ReLU activation and dropout with $p=0.2$, containing around 5 million parameters. The implementation supports flexible configurations, enabling the construction and evaluation of more complex deep classifier heads for other use cases.

Both LoRA adapters and the X-LoRA scaling head are trained using next-token prediction with cross-entropy loss (we compute the probability that the model assigns to the correct next token in a sequence based on the preceding tokens; during training, the objective is to maximize the likelihood of the correct next token as indicated in the training dataset). We apply token shifting, where the input sequence is shifted by one token to form the target sequence that the model aims to predict. For instance, if the input sequence is ``Proteins are fundamental to", the corresponding target sequence for training would be ``are fundamental to biology". Consequently, the model endeavors to predict the subsequent token in the sequence.

As illustrated in Figure~\ref{fig:hierarchical_concept}, during training of the X-LoRA scaling head, only the scaling head parameters are updated while all other model parameters remain fixed. The original Zephyr tokenizer and prompt template are employed.

\subsection{X-LoRA-Gemma model} The X-LoRA-Gemma model is constructed similarly to the previously described X-LoRA model but is based on the Gemma-7B-it model~\cite{Google/gemma-7b-itFace}, utilizing four adapters (refer to the main text for further details). These four adapters are trained on datasets covering mechanics and materials, protein mechanics, bioinspired materials, as well as an additional quantum-mechanics-based molecular properties dataset QM9 (refer to Table~\ref{tab:table_QM9})~\cite{Ruddigkeit2012EnumerationGDB-17,Ramakrishnan2015ElectronicSpace}.

Each LoRA adapter has rank $r=16$ with scaling factor $\alpha=16$. The adapters are trained targeting seven modules: \texttt{q\_proj}, \texttt{o\_proj}, \texttt{k\_proj}, \texttt{v\_proj}, \texttt{gate\_proj}, \texttt{up\_proj}, and \texttt{down\_proj}, using a dropout rate of 0.05. Given the four adapters and seven target modules, and considering the 28 layers in the Gemma-7B base model, a total of $28 \times 28 = 784$ values of $\lambda_i$ are calculated for each token.

The first two adapters (bioinspired and mechanics of materials) included in this model are trained using question-answer pairs provided by the respective datasets (see Table~\ref{tab:table_loradata}).

Protein mechanics tasks are trained similarly to the previous section, employing bidirectional forward and inverse instruction sets. Additionally, an adapter is trained on the QM9 dataset with the following two tasks:

\begin{quote}
\tiny {\texttt{CalculateMolecularProperties<CC(C)N1CC1C=O> [0.098,0.358,0.581,0.395,0.309,0.330,0.570,0.649,0.519,0.519,0.519,0.518]\\
GenerateMolecularProperties<0.098,0.358,0.581,0.395,0.309,0.330,0.570,0.649,0.519,0.519,0.519,0.518> [CC(C)N1CC1C=O]\\ 
...
}}
\end{quote}

These correspond to two tasks in total (predicting molecular properties, the forward task, and the inverse task of generating a molecule with specified target properties). All 12 properties included in the QM9 dataset are either predicted or designed for, respectively (refer to Table~\ref{tab:table_QM9}).  
The model incorporates four adapters in total.

The X-LoRA scaling head for this model consists of three layers, with the intermediate layer dimension being $3072\times {FF}_{\rm fac}=12,288$, where ${FF}_{\rm fac}=4$ is the feed-forward multiplier. Linear layers with ReLU activation functions and a dropout rate of $p=0.2$ are employed, resulting in approximately 47 million parameters.

The X-LoRA-Gemma scaling head is trained on about 17,000 samples. These samples were selected as a subset of the training dataset used for adapter training; notably, although the X-LoRA-Gemma model had only four adapters, the samples were drawn from the full training dataset of the previously described model, with the addition of samples from the QM9 dataset. This approach was taken to test whether the X-LoRA model could acquire further capabilities during the scaling head training phase. Given that the loss converges well to 0.58 by the end of training, alongside strong performance on tasks the model was not explicitly trained on (e.g., chemistry, as demonstrated by the molecule design example in the paper), we consider this approach to generally produce favorable outcomes.

We employed a \texttt{paged\_adamw\_8bit} optimizer with a gradient clipping threshold of 0.3, a learning rate of $1\times 10^{-4}$ incorporating warmup, and four steps of gradient accumulation, using a batch size of one. The X-LoRA-Gemma model was trained for approximately 62,000 steps.

The original Gemma tokenizer and prompt template were utilized.

\subsection{Adversarial agentic modeling} We realize an adversarial agentic approach by creating two X-LoRA agents. One agent specializes in formulating questions. It poses the initial question and continuously generates further questions aimed at uncovering challenges and weaknesses in the answers. The other agent handles responding to these queries, providing ongoing answers. This approach is implemented through \{Guidance\}~\cite{Guidance-ai/guidance:Models.}.

Regarding the example connecting music and mechanics presented in the text, the roles of the two agents are defined as follows. Firstly, the physicist and question asker:

\begin{quote}
\small\texttt{You are a critical physicist engaged in a discussion from a physics standpoint. Keep your responses concise and always challenge assertions provocatively.

As a creative thinker, you incorporate ideas from other domains and push conceptual boundaries.}
\end{quote}

Secondly, a philosopher (in certain cases, this role is adapted to represent a biology expert):

\begin{quote}
\small\texttt{You are a philosopher well-versed in biology, chemistry, and mathematics participating in a discussion.

Maintain your answers brief yet accurate and imaginative. You generate excellent ideas and novel lines of reasoning, always grounded in logic.}
\end{quote}

The question asker is specifically directed to reflect on the current question and previous answers and to craft new probing questions. This is achieved through a prompt structured as follows:

\begin{quote}
\small\texttt{Examine the following question and answer.

\#\#\# Question: <question>

\#\#\# Response: <response>

\#\#\# Instruction: Provide a SINGLE follow-up question that critically examines the response.

Do NOT provide an answer or any commentary yet.

The single question is: }
\end{quote}

After the dialogue has progressed through $N$ turns, the model is assigned to 1) summarize the main insights discussed, 2) enumerate the key insights in bullet form, and 3) pinpoint the most significant takeaway.

The raw conversation text serves as the foundation for subsequent analysis through graph construction.

\subsection{Knowledge graph generation} We employ Zephyr-7B-$\beta$ to extract triplets from the text, following the methodology described in~\cite{Rahulnyk/knowledge_graph:QnA} with enhancements inspired by the Llama Index graph construction method~\cite{Liu_LlamaIndex_2022}. The graphs are visualized using NetworX\cite{Networkx/networkx:Python} and Pyvis~\cite{WestHealth/pyvis:Graphs.}. A principal instruction given is, adhering to the technique recommended in~\cite{Liu_LlamaIndex_2022}:

\begin{quote}
\small\texttt{Your task is to extract the ontology of terms mentioned within the provided context. These terms should embody the central concepts relevant to the context.

Format your output as a list of JSON objects. Each list element contains a pair of terms and the relation connecting them, as illustrated by: [...] }
\end{quote}

We direct the creation of triplets consisting of nodes and edge attributes as follows, following the procedure in \cite{Liu_LlamaIndex_2022}:

\begin{quote}
\small\texttt{\begin{itemize}
            \item "node\_1": "A concept identified in the extracted ontology"
            \item "node\_2": "A related concept identified in the extracted ontology"
            \item "edge": "A concise description of the relationship between node\_1 and node\_2"
        \end{itemize}
        }
\end{quote}

We also supply an example of ontological analysis to the model, for instance:

\begin{quote}
\small\texttt{
       Context: ```Spider silk is a strong natural fiber used to catch prey in a web.``` }
\end{quote}

Subsequently, example triplets are provided, such as:

\begin{quote}
\small\texttt{\begin{itemize}
            \item "node\_1": "spider silk"
            \item "node\_2": "fiber"
            \item "edge": "is"
        \end{itemize}
        }
\end{quote}

The generated graph data is then compiled and clustered into communities applying the Girvan-Newman algorithm~\cite{Girvan2002CommunityNetworks}.

\subsection{Visualization of molecular structures}

PyMOL~\cite{Schrodinger2015The1.8} is utilized to visualize and investigate the predicted protein structures. Various scripts and functions in this program enable the identification of specific protein features such as secondary structures, hydrophobic and hydrophilic regions, disulfide bonds, and hydrogen bonds.

\section*{Data Availability Statement} The code and data supporting the outcomes of this research are publicly accessible at \url{https://github.com/EricLBuehler/xlora} and \url{https://huggingface.co/lamm-mit/x-lora}. Additional datasets employed in this work are cited (see particularly Table~\ref{tab:table_loradata}).

\end{document}